% Hoe dit boek te lezen: a \subchapter-level (H2) heading grouped under the
% single "Inleiding" chapter-level heading in main.tex.
\subchapter{Hoe dit boek te lezen}

\lettrine{H}{et} boek is ingedeeld per gang, van voorgerecht tot dessert. De
hoofdgerechten zijn verder opgesplitst naar type: pasta en risotto, quiches,
stamppot, wereldkeukens. Voorin zit een inhoudsopgave om per categorie te
bladeren, achterin een register om op gerecht of ingrediënt op te zoeken.

Na dit hoofdstuk staan nog een paar korte, praktische hoofdstukken voordat de
recepten beginnen. De seizoenskalender laat zien wanneer onze groente en
fruit op hun best zijn. De basisvoorraad is wat ik zelf altijd in de kast heb
staan. Portiegrootte legt uit met hoeveel personen dit boek standaard
rekent. En kleine maten geeft een indicatie in gram voor vage hoeveelheden
als een snufje of een scheutje.

Elk recept begint met een kickerregel die het gerecht in het kort typeert:
het soort gang, of het vlees, vis, vegetarisch of vegan is, en de keuken.
Daaronder staat de bereidingstijd, en als dat afwijkt van de standaard vijf
porties ook het aantal porties. Bij de ingrediënten staat de hoeveelheid
links van de naam; een \textcolor{terracotta}{\raisebox{0.15ex}{$\scriptstyle\blacklozenge$}}
ervoor betekent dat er geen vaste hoeveelheid bij hoort, zoals een snufje
zout of een scheutje olie. In de kantlijn duikt af en toe een Tip van Ben
op, in het handschrift, met iets wat handig is om te weten voordat je
begint. Lees een recept voor je begint het handigst even helemaal door,
zodat je niet halverwege ontdekt dat er nog iets een uur moet marineren.

De kicker is bedoeld om in één oogopslag te zien wat voor gerecht het is,
niet om allergieën of diëten af te dekken: glutenvrij, lactosevrij of
notenvrij staat er nergens bij. Heb je daar rekening mee te houden, check
dan altijd zelf de ingrediëntenlijst.

Staat er in een recept gewoon olijfolie, dan bedoel ik daarmee altijd extra
vergine olijfolie, ook om in te bakken. Een goede olijfolie is voor mij een
basisingrediënt, niet iets wat je bewaart voor achteraf. Sommige mensen
zijn bang dat extra vergine olijfolie bij verhitting zijn kwaliteit
verliest, maar dat blijkt niet te kloppen: onderzoek laat zien dat hij bij
hoge temperaturen juist stabieler blijft dan de meeste andere
kookoliën~\cite{dealzaa-olijfolie-verhitting}.

Bij de brood-, gebak- en koekjesrecepten staat achter water, zout, vet en
suiker ook een percentage, bijvoorbeeld water (75\%). Dat is het
bakkerspercentage: het gewicht van dat ingrediënt uitgedrukt als percentage
van het bloemgewicht. Handig als je een recept wilt opschalen, of wilt
vergelijken met een ander recept.

Dit boek staat trouwens ook online. Daar kun je de recepten doorzoeken en
gewoon in de browser lezen\cite{kookboek-online}.
